% LaTeX resume using res.cls
\documentclass[line,margin]{res} 
\usepackage{helvetica} % uses helvetica postscript font (download helvetica.sty)
\usepackage{hyperref,lmodern}
\hyphenation{Altium Designer}
\newcommand{\CC}{C\nolinebreak\hspace{-.05em}\raisebox{.4ex}{\tiny\bf +}\nolinebreak\hspace{-.10em}\raisebox{.4ex}{\tiny\bf +}}

\begin{document}

\name{Maya Winland-Gaetz}
\address{203-1870 McKenzie Avenue, Victoria BC V8N 4X3}
\address{\href{mailto:mwinlandgaetz@gmail.com}{mwinlandgaetz@gmail.com}\\
(289) 689-0427}

\begin{resume}
\section{Education}
    {\textbf{Bachelor of Engineering (BEng),}} In Progress, GPA 8.58/9 \hfill Expected Sept 2028\\
    Electrical Engineering\\
    \textbf{University of Victoria}, Victoria, BC, Canada \\\\
    \textbf{PhD Student,} (Incomplete), GPA 11.0/12   \hfill September 2013 - May 2019\\
    \textbf{McMaster University}, Hamilton, ON, Canada \hfill 

    {\textbf{Bachelor of Science (BSc),}} GPA 3.57/4 \hfill September 2008 - April 2013\\
    Specialist Environmental Toxicology, Maj. Ecology \& Evolutionary Biology\\
    \textbf{University of Toronto}, Toronto, ON, Canada \\\\


\section{Competencies}
    \textbf{Technical}
    \begin{itemize} \itemsep -2pt
        \item Drafted, simulated, and am currently in fabrication for several electronics projects using \textbf{KiCAD}, \textbf{LTSpice}, and \textbf{Altium Designer}.
        \item Developed several applications in \textbf{\CC} for personal use.
        \item 90WPM typing speed with a keen eye for typographical and grammatical errors, using both \textbf{Microsoft Office} and \textbf{LaTeX}.
        \item Prepared and analyzed sensitive geochemical samples in a laboratory environment, implementing state-of-the-art instruments while observing rigorous safety protocols.
    \end{itemize}
    \textbf{Core}
    \begin{itemize} \itemsep -2pt
        \item Maintained efficiency while under pressure during the annual Starbucks Holiday Launch, providing a high standard of service for eager early-morning customers.
        \item Communicated clearly and fluently in presentations and public speaking, earning an award at an academic conference for a 15 minute talk on geochemical sampler design.
    \end{itemize}

\section{Professional Experience}
    \textbf{Electrical Engineering Co-op} \hfill September 2025 - December 2025 \\
    \textbf{Ocean Networks Canada}, Sidney, BC, Canada
    \begin{itemize}  \itemsep -2pt
        \item Developed a data capture, storage, and graphing application for the Spanish Antarctica project using Python and SQLite, which directly contributed to the detection and resolution of several mission-critical bugs prior to station deployment.
        \item Compiled an exhaustive feasibility study for an upgrade to the power distribution system for the Pacific Ocean Neutrino Experiment, including extensive component and connector research, circuit design and multi-board modelling in Altium Designer, and transient analysis using LTSpice.
        \item Reverse-engineered the SPI communication protocol for a deprecated Alcatel fiber-optic node using a mixed-signal oscilloscope/digital logic analyzer, and programmed an Arduino UNO in C++ to intercept SPI communications and issue commands to the tunable laser to ensure lasting functionality of the nodes in the ONC infrastructure.
        \item Prepared comprehensive reports and documentation on all the above projects, to ensure future co-ops/employees can build upon my work.
    \end{itemize} 
    \textbf{Teaching Assistant} \hfill September 2013 - April 2018 \\
    \textbf{McMaster University}, Hamilton, ON, Canada
    \begin{itemize}  \itemsep -2pt
        \item Planned and conducted lab sessions and tutorials, using previous experience and WHMIS training in chemistry lab safety and analytical protocols, to minimize accidents, control spills and ensure a safe learning environment.
        \item Graded and annotated student papers and projects on 7-day deadlines, using rigorous organization and uniform standards, to ensure course leaders could maintain a strict schedule.
    \end{itemize} 

\section{Projects} 
\textbf{University of Victoria Satellite Design Club} \hfill January 2024 - Present
\begin{itemize} \itemsep -2pt
    \item Researched satellite radiocommunication design principles and protocols to guide the development of the telemetry systems for future club projects.
    \item Responsible for the design and assembly of the Telemetry, Tracking \& Command (TT\&C) boards and evaluation modules for the SLIP payload project, a partnership with UVic Rocketry.
    \item Hand-soldered 7mm QFN and 0603/0402-sized components for the board using a hot air rework station and manual soldering iron to assemble the first iteration of transceiver evaluation boards for testing.
\end{itemize}
\textbf{First Year Engineering Design Project} \hfill January 2024 - April 2024
\begin{itemize} \itemsep -2pt
    \item Conceived and constructed a physical prototype and associated webapp for a smart shower heater as part of a class project. The product could sense the proximity of a user and automatically adjust the water temperature using a hypothetical inline water heater to a desired temperature, configurable through a WiFi-enabled app.
    \item Went beyond the expectations of the course by building transistor amplifiers to facilitate proximity detection with inexpensive IR sensors, and improved the TA-provided netcode to ensure a stable connection.
    \item Worked in a group of 3, wherein the circuit design and netcode improvements were my chief responsibilities.\\
    \textit{Demonstration video available upon request. }
\end{itemize}
\textbf{Personal} \hfill April 2020-Present
\begin{itemize} \itemsep -2pt
    \item Currently building a trigger signal sequencer as part of a larger project to construct a music synthesizer, utilizing RP2040 hardware timers and digital logic ICs to output a 4-channel beat structure with a custom LED matrix.
    %\item Researched components for and prototyped auxiliary modules, including a headphone output amplifier and switch-mode power supply for the overall project.
    \item Programmed and assembled a climate and light level sensor for personal use, incorporating an AtMega328P microcontroller and a scrolling 16x8 LED display.
    \item Programmed an app in C++ to view a directory of life drawing references at regular intervals, to aid in life drawing study during the pandemic.
    %\item Implemented a ray tracer/ray marcher algorithm in C++ using the SFML graphics library, rendering simple 3-dimensional objects to a static output.
\end{itemize}




\section{Awards and Certifications}
\begin{itemize} \itemsep -2pt
    \item Amateur Radio Operator Certification\hfill December 2024\\ Advanced, Basic with Honours (VA7MWX)
    \item Best Student Oral Presentation, worth \$750  {\sl (Sudbury Mining and Environment Int'l Conference)} \hfill June 2015
    \item WHMIS \hfill Yearly Recertification 2008-2022
    % \begin{itemize} \itemsep -2pt
    %     \item Other safety certifications as needed (Hazardous materials disposal and compressed gases)
    % \end{itemize}
\end{itemize}
\section{Other Interests}
\begin{itemize} \itemsep -2pt
    \item Bicycling, with the goal of getting from the University of Victoria to Leechtown and back over a two day trip.
    \item Classical piano, playing Beethoven’s later Sonatas in particular.
\end{itemize}
\textbf{References Available Upon Request}
\end{resume}
\end{document}







